\begin{proposition}[Cauchy criterion for convergence of series of real numbers] \label{proposition:cauchy_criterion_for_convergence_of_series_of_real_numbers}
    \TODO{AI generated, verify}
    Let $(a_n)_{n=1}^\infty$ be a sequence in $\mathbb{R}$\CrefIfExists{definition:field_of_real_numbers} and let $S_N = \sum_{n=1}^N a_n$ be the sequence of partial sums.

    The series $\sum_{n=1}^\infty a_n$ \CrefAndHyperrefIfExist{definition:convergence_and_limit_of_sequence_in_extended_metric_space}{converges} if and only if for every $\varepsilon > 0$ there exists $N \in \mathbb{N}$ such that for all $M \geq N$,
    $$\left| \sum_{n=N+1}^M a_n \right| < \varepsilon.$$
    Equivalently, the series converges if and only if the sequence of partial sums $(S_N)$ is a \CrefAndHyperrefIfExist{definition:cauchy_sequence_in_an_extended_metric_space}{Cauchy sequence} in $\mathbb{R}$\CrefIfExists{definition:field_of_real_numbers}.
\end{proposition}
\begin{proof}
    \TODO{AI generated, verify}
    The series $\sum_{n=1}^\infty a_n$ converges if and only if the sequence of partial sums $(S_N)$ converges in $\mathbb{R}$. Since $\mathbb{R}$ is complete (every Cauchy sequence converges), a sequence in $\mathbb{R}$ converges if and only if it is a Cauchy sequence.

    The sequence $(S_N)$ is Cauchy means that for every $\varepsilon > 0$ there exists $N \in \mathbb{N}$ such that for all $M, K \geq N$,
    $$|S_M - S_K| < \varepsilon.$$

    Without loss of generality assume $M \geq K$. Then
    $$S_M - S_K = \sum_{n=1}^M a_n - \sum_{n=1}^K a_n = \sum_{n=K+1}^M a_n.$$

    Therefore the Cauchy condition on $(S_N)$ is equivalent to: for every $\varepsilon > 0$ there exists $N \in \mathbb{N}$ such that for all $M \geq N$,
    $$\left| \sum_{n=N+1}^M a_n \right| < \varepsilon.$$

    This completes the proof.
\end{proof}
